\chapter{La boucle QUBO bi-objectif}\label{ch:boucle}
\iterbanner{V1}

\section{L'algorithme}
La note d'intention décrit une boucle itérative : simuler quelques profils, ajuster un modèle quadratique, en tirer un QUBO, le résoudre par QAOA ou recuit, simuler les meilleurs candidats, recommencer. Le cours quantique l'a appliquée à la seule traînée. Ici, on l'étend aux deux objectifs et aux contraintes.

\begin{figure}[htbp]
\centering
\begin{tikzpicture}[node distance=0.5cm and 0.55cm]
\node[flow] (init) {10 profils\\ tirés au hasard};
\node[flow,right=of init] (fit) {ajuster $\widehat{C_D}$, $\widehat{\mathrm{SER}}$\\ (quadratiques)\\ et $\widehat{C_L}$ (affine)};
\node[flow,right=of fit] (qubo) {pour chaque $w$ :\\ QUBO à 13 qubits\\ objectifs + contraintes};
\node[flow,right=of qubo] (solve) {solveur :\\ exact, recuit\\ ou QAOA};
\node[flow,below=1.1cm of qubo] (sim) {simuler les meilleurs\\ profils \textbf{non encore simulés}};
\node[flow,below=1.1cm of fit] (arch) {archive : profils simulés\\ front = non dominés de l'archive};
\draw[arr] (init) -- (fit); \draw[arr] (fit) -- (qubo); \draw[arr] (qubo) -- (solve);
\draw[arr] (solve) |- (sim); \draw[arr] (sim) -- (arch); \draw[arr] (arch) -- (fit);
\end{tikzpicture}
\caption{La boucle bi-objectif. Le front rendu n'est formé que de profils réellement simulés, comme l'exige la note d'intention.}\label{fig:boucle}
\end{figure}

\section{Étape 1 : les modèles de substitution}
On ajuste par moindres carrés, sur les profils faisables déjà simulés, deux polynômes quadratiques en variables normalisées $z_i=2i_i/7-1$ (10 coefficients chacun, « famille A » du cours aérodynamique) :
\[
\widehat{C_D}(z),\ \widehat{\mathrm{SER}}(z)=c_0+\sum_ic_iz_i+\sum_{i\le l}c_{il}z_iz_l .
\]
Pour la contrainte de portance, on ajuste un modèle \emph{affine} $\widehat{C_L}$ sur les profils valides : $C_L$ est quasi affine en $(m,\alpha)$, et une contrainte doit rester de degré~1 pour que sa pénalité soit quadratique.

Une petite régularisation (ridge, $10^{-6}$) évite les systèmes mal conditionnés quand peu de points sont disponibles. Avec 23 profils faisables, sur le scénario S2, le modèle de SER atteint $R^2=0{,}996$ sur toute la grille, celui de traînée $R^2=0{,}86$ : la SER, lisse, est plus facile à modéliser que la traînée et ses bascules de transition.

\section{Étape 2 : un QUBO par poids}
\paragraph{L'objectif.} Un QUBO n'a qu'une fonction de coût : on utilise une somme pondérée des deux modèles, normalisés par la plage des valeurs déjà observées,
\[
f_w=w\,\frac{\widehat{C_D}}{\Delta C_D}+(1-w)\,\frac{\widehat{\mathrm{SER}}}{\Delta\mathrm{SER}},\qquad w\in\{0;\ 0{,}25;\ 0{,}5;\ 0{,}75;\ 1\}.
\]
Comme la substitution $z=a_0+A\,\vect b$ est linéaire, on combine d'abord les coefficients, puis on développe en bits une seule fois (\code{surrogate\_to\_qubo}). On obtient 9 qubits.

\paragraph{La contrainte d'épaisseur : un qubit auxiliaire.} $t\ge8$~\% interdit l'indice $i_t=0$, c'est-à-dire les trois bits d'épaisseur nuls. La pénalité naturelle, $P(1-b_0)(1-b_1)(1-b_2)$, est \emph{cubique} : ce n'est pas un QUBO (c'est le piège du degré vu dans le cours quantique). On la rend quadratique par la \textbf{réduction de Rosenberg}. On introduit un qubit $y$ censé valoir $x_1x_2$, avec $x_1=1-b_1$ et $x_2=1-b_2$, et la pénalité
\[
P\,(1-b_0)\,y+2P\big[x_1x_2-2x_1y-2x_2y+3y\big].
\]
Le crochet vaut 0 si $y=x_1x_2$ et au moins 1 sinon. Avec un poids $2P>P$, le minimum sur $y$ est atteint pour $y=x_1x_2$, et on retrouve exactement la pénalité cubique. Le code vérifie l'égalité sur les 8 combinaisons par énumération.

\paragraph{La contrainte de portance : des qubits d'écart.} $\widehat{C_L}(\vect b)\ge0{,}3$ est une inégalité. On la transforme en égalité $\widehat{C_L}(\vect b)-0{,}3-s=0$ avec une variable d'écart $s\ge0$ codée sur 3 qubits, et on pénalise le carré :
\[
P_{C_L}\big(\widehat{C_L}(\vect b)-0{,}3-s\big)^2,\qquad s=\frac{s_{\max}}{7}\sum_{j=0}^{2}2^js_j .
\]
Comme $\widehat{C_L}$ est affine, le carré est bien quadratique. Contrepartie : $s$ ne prend que 8 valeurs, et le résidu minimal n'est pas toujours nul. La pénalité ajoute donc un petit biais aux profils faisables, d'autant plus faible que le pas de $s$ est fin.

\begin{codebox}[ : \code{optim/qubo\_loop.py}]
\begin{lstlisting}
def build_qubo(coef_cd, coef_ser, coef_cl, w, scale_cd, scale_ser, cons, P=3.0, n_slack=3):
    coef = w * coef_cd / scale_cd + (1 - w) * coef_ser / scale_ser
    Q, c = surrogate_to_qubo(coef, BITS, LEVELS)              # 9 qubits
    Q = grow(Q, 1 + n_slack)                                   # + y + 3 bits d'ecart
    Q, c = penalty_index_nonzero(Q, c, bits=(6, 7, 8), y=9, P=P)          # t >= 8 %
    const, lin = lin_bits(coef_cl)                             # C_L affine -> forme affine en bits
    Q, c = penalty_linear_geq(Q, c, lin, const, 0.3, [10, 11, 12], s_max, 25 * P)   # C_L >= 0.3
    return Q, c, 13
\end{lstlisting}
\end{codebox}

\begin{figure}[htbp]
\centering
% figure générée par tools/make_figures.py : matrice_Q
\begin{tikzpicture}
\begin{axis}[width=7.4cm,height=7.4cm,enlargelimits=false,axis on top,y dir=reverse,
  xtick={0,...,12},ytick={0,...,12},xticklabels={$m_0$,$m_1$,$m_2$,$p_0$,$p_1$,$p_2$,$t_0$,$t_1$,$t_2$,$y$,$s_0$,$s_1$,$s_2$},yticklabels={$m_0$,$m_1$,$m_2$,$p_0$,$p_1$,$p_2$,$t_0$,$t_1$,$t_2$,$y$,$s_0$,$s_1$,$s_2$},
  tick label style={font=\scriptsize},colorbar,colorbar style={width=2.5mm,tick label style={font=\scriptsize}},
  colormap={div}{color=(corange) color=(white) color=(cblue)},point meta min=-1,point meta max=1]
\addplot[matrix plot*,mesh/cols=13,point meta=explicit] coordinates {(0,0) [-0.0524] (1,0) [0.1171] (2,0) [0.2342] (3,0) [0.0085] (4,0) [0.0170] (5,0) [0.0340] (6,0) [0.0072] (7,0) [0.0144] (8,0) [0.0289] (9,0) [0.0000] (10,0) [-0.0625] (11,0) [-0.1250] (12,0) [-0.2500] (0,1) [0.1171] (1,1) [-0.0462] (2,1) [0.4684] (3,1) [0.0170] (4,1) [0.0340] (5,1) [0.0681] (6,1) [0.0144] (7,1) [0.0289] (8,1) [0.0577] (9,1) [0.0000] (10,1) [-0.1250] (11,1) [-0.2500] (12,1) [-0.5000] (0,2) [0.2342] (1,2) [0.4684] (2,2) [0.1419] (3,2) [0.0340] (4,2) [0.0681] (5,2) [0.1362] (6,2) [0.0289] (7,2) [0.0577] (8,2) [0.1154] (9,2) [0.0000] (10,2) [-0.2500] (11,2) [-0.5000] (12,2) [-1.0000] (0,3) [0.0085] (1,3) [0.0170] (2,3) [0.0340] (3,3) [-0.0116] (4,3) [0.0036] (5,3) [0.0071] (6,3) [0.0010] (7,3) [0.0020] (8,3) [0.0040] (9,3) [0.0000] (10,3) [-0.0094] (11,3) [-0.0188] (12,3) [-0.0376] (0,4) [0.0170] (1,4) [0.0340] (2,4) [0.0681] (3,4) [0.0036] (4,4) [-0.0213] (5,4) [0.0143] (6,4) [0.0020] (7,4) [0.0040] (8,4) [0.0080] (9,4) [0.0000] (10,4) [-0.0188] (11,4) [-0.0376] (12,4) [-0.0752] (0,5) [0.0340] (1,5) [0.0681] (2,5) [0.1362] (3,5) [0.0071] (4,5) [0.0143] (5,5) [-0.0355] (6,5) [0.0040] (7,5) [0.0080] (8,5) [0.0159] (9,5) [0.0000] (10,5) [-0.0376] (11,5) [-0.0752] (12,5) [-0.1504] (0,6) [0.0072] (1,6) [0.0144] (2,6) [0.0289] (3,6) [0.0010] (4,6) [0.0020] (5,6) [0.0040] (6,6) [-0.0059] (7,6) [0.0011] (8,6) [0.0021] (9,6) [-0.0733] (10,6) [-0.0076] (11,6) [-0.0152] (12,6) [-0.0303] (0,7) [0.0144] (1,7) [0.0289] (2,7) [0.0577] (3,7) [0.0020] (4,7) [0.0040] (5,7) [0.0080] (6,7) [0.0011] (7,7) [-0.1579] (8,7) [0.1509] (9,7) [0.2933] (10,7) [-0.0152] (11,7) [-0.0303] (12,7) [-0.0606] (0,8) [0.0289] (1,8) [0.0577] (2,8) [0.1154] (3,8) [0.0040] (4,8) [0.0080] (5,8) [0.0159] (6,8) [0.0021] (7,8) [0.1509] (8,8) [-0.1671] (9,8) [0.2933] (10,8) [-0.0303] (11,8) [-0.0606] (12,8) [-0.1213] (0,9) [0.0000] (1,9) [0.0000] (2,9) [0.0000] (3,9) [0.0000] (4,9) [0.0000] (5,9) [0.0000] (6,9) [-0.0733] (7,9) [0.2933] (8,9) [0.2933] (9,9) [-0.0733] (10,9) [0.0000] (11,9) [0.0000] (12,9) [0.0000] (0,10) [-0.0625] (1,10) [-0.1250] (2,10) [-0.2500] (3,10) [-0.0094] (4,10) [-0.0188] (5,10) [-0.0376] (6,10) [-0.0076] (7,10) [-0.0152] (8,10) [-0.0303] (9,10) [0.0000] (10,10) [0.1207] (11,10) [0.1340] (12,10) [0.2681] (0,11) [-0.1250] (1,11) [-0.2500] (2,11) [-0.5000] (3,11) [-0.0188] (4,11) [-0.0376] (5,11) [-0.0752] (6,11) [-0.0152] (7,11) [-0.0303] (8,11) [-0.0606] (9,11) [0.0000] (10,11) [0.1340] (11,11) [0.3085] (12,11) [0.5361] (0,12) [-0.2500] (1,12) [-0.5000] (2,12) [-1.0000] (3,12) [-0.0376] (4,12) [-0.0752] (5,12) [-0.1504] (6,12) [-0.0303] (7,12) [-0.0606] (8,12) [-0.1213] (9,12) [0.0000] (10,12) [0.2681] (11,12) [0.5361] (12,12) [0.8850]};
\end{axis}
\end{tikzpicture}

\caption{Un QUBO de la boucle (scénario S2, $w=0{,}5$, après 30 simulations), normalisé par son plus grand coefficient. Le bloc $9\times9$ en haut à gauche vient des modèles et du couplage avec la contrainte de portance ; la colonne $y$ et le bloc $s$ viennent des pénalités. Au total, 69 couplages non nuls, soit 138 CNOT par couche de QAOA.}\label{fig:matQ}
\end{figure}

\section{Étape 3 : résoudre et choisir les candidats}
Trois solveurs sont disponibles, de même interface : ils renvoient une liste de profils classés par énergie croissante.
\begin{itemize}
\item \textbf{exact} : énumération des $2^{13}=8\,192$ états. Il sert de référence, en isolant l'effet du modèle de substitution de celui du solveur ;
\item \textbf{recuit simulé} : 40 lectures, 150 balayages, comme le \code{SimulatedAnnealingSampler} d'Ocean ;
\item \textbf{QAOA} : profondeur $p=2$, énergie normalisée, 3 départs pour COBYLA, puis 512 tirs ; on trie les chaînes distinctes obtenues par leur énergie recalculée classiquement.
\end{itemize}
Pour chaque poids, on retient le meilleur profil \emph{qui n'a pas encore été simulé}. Ce détail est essentiel : sans lui, la boucle proposerait indéfiniment le même optimum du modèle, déjà connu.

\paragraph{Les points non supportés.} La somme pondérée ne vise que les 4 ou 5 points supportés du front (chapitre~\ref{ch:exact}). La boucle trouve pourtant les autres, pour deux raisons : l'exclusion des profils déjà simulés pousse vers les deuxièmes, troisièmes meilleurs, qui sont souvent des voisins du front ; et l'imperfection du modèle quadratique déplace les minima d'une itération à l'autre. C'est une exploration involontaire mais efficace.

\section{Étape 4 : simuler, archiver, recommencer}
Les candidats sont simulés (le compteur avance), ajoutés à l'archive, et le front courant est recalculé sur les seuls profils simulés. On s'arrête quand le budget est atteint, ou, pour la mesure de performance, dès que tout le front exact a été simulé.

\exercices
\begin{exercice}
Vérifier la réduction de Rosenberg : pour $x_1=x_2=1$, comparer la pénalité obtenue pour $y=0$ et $y=1$, et montrer que le minimum vaut $P(1-b_0)$. Faire de même pour $x_1=1$, $x_2=0$.
\end{exercice}
\begin{exercice}
Combien de qubits faudrait-il si l'on codait la contrainte $i_t\ge1$ par une variable d'écart, comme celle de portance, au lieu du qubit auxiliaire ? Quel avantage a la réduction de Rosenberg ici ?
\end{exercice}

\begin{pistes}
\begin{itemize}
\item Remplacer la somme pondérée par une \textbf{$\varepsilon$-contrainte} $\min\widehat{C_D}$ sous $\widehat{\mathrm{SER}}\le\varepsilon$ : elle vise aussi les points non supportés, au prix de qubits d'écart supplémentaires.
\item Pondérer les moindres carrés vers les meilleurs profils (modèle \emph{local}, région de confiance), pour que le modèle soit précis là où se trouve le front.
\item Choisir les candidats par un critère d'\emph{amélioration attendue de l'hypervolume}, comme en optimisation bayésienne multi-objectif (ParEGO, EHVI) : c'est le standard actuel des méthodes à modèle de substitution, auquel la boucle QUBO devrait être comparée.
\item Remplacer la variable d'écart de $C_L$, qui introduit un biais, par une réduction exacte quand la contrainte se ramène à une borne sur une variable (ici, $C_L\ge0{,}3$ équivaut à $m\ge1$~\% sur la grille).
\end{itemize}
\end{pistes}
